\begin{lemma}[An integral dominated by a product of two $L^2$ functions]
Fix $n \in \mathbb{N}_0$ and write $\lambda$ for Lebesgue measure on
$\mathbb{R}^n$.  Let $K \subseteq \mathbb{R}^n$ and let
$\lambda_K$ denote the restriction of $\lambda$ to $K$, so that for a function
$f$ one writes $\|f\|_{L^2(K)} = \|f\|_{L^2(\lambda_K)}$ and
$\int_K f \, d\lambda$ for the integral against $\lambda_K$.  Let
$h, a, b : \mathbb{R}^n \to \mathbb{R}$ and let $C \in \mathbb{R}$ with
$C \geq 0$.  Assume
\begin{enumerate}
  \item $a \in L^2(\lambda_K)$ and $b \in L^2(\lambda_K)$;
  \item $|h(x)| \leq C\,\bigl(|a(x)|\,|b(x)|\bigr)$ for every
    $x \in \mathbb{R}^n$.
\end{enumerate}
Then
\[
  \Bigl\| \int_K h \, d\lambda \Bigr\|
    \leq C\,\|a\|_{L^2(K)}\,\|b\|_{L^2(K)} .
\]

Two points about the hypotheses, both of which are exactly as stated and not
abbreviations of something stronger.  First, every measure here is Lebesgue:
$\lambda_K$ is Lebesgue measure restricted to $K$ and the $L^2(K)$ norms are
taken against it, with no Gaussian weight anywhere, even though the lemma is
used inside the tablet's Gaussian development.  Second, nothing at all is
assumed about $h$ beyond the pointwise bound (2) --- no integrability and no
measurability --- and no such hypothesis is needed: if $h$ fails to be
integrable then the Bochner integral $\int_K h \, d\lambda$ is by convention
$0$, while the right-hand side is nonnegative because $C \geq 0$ and norms are
nonnegative, so the inequality holds a fortiori.  Correspondingly the lemma
cites no other tablet node; it rests only on the shared background measure
theory.
\end{lemma}
\begin{proof}
Step 1 (Cauchy--Schwarz on $K$).  We first record that for $a, b \in
L^2(\lambda_K)$,
\[
  \int_K |a|\,|b| \, d\lambda
    \leq \|a\|_{L^2(K)}\,\|b\|_{L^2(K)} . \tag{CS}
\]
Indeed, the product $ab$ lies in $L^1(\lambda_K)$, by the H\"older pairing of
the exponents $2$ and $2$ into $1$, so it is $\lambda_K$-integrable.  Since
$\bigl\| a(x)b(x) \bigr\| = |a(x)|\,|b(x)|$ at every $x$, the integral on the
left of (CS) is the $L^1(\lambda_K)$ norm of $ab$, read as a real number; and
the H\"older inequality for the exponent pairing $1/2 + 1/2 = 1$ bounds that
$L^1$ norm by the product of the two $L^2$ norms.  Both sides are finite --- the
right side because $a$ and $b$ are in $L^2(\lambda_K)$, the left because of the
bound itself --- so passing from the inequality between extended-real norms to
the inequality (CS) between their real values is legitimate, the monotonicity of
$t \mapsto t$ on $[0,\infty]$ restricted to finite values being all that is
used.

Step 2 (the integrand of the majorant is integrable).  The function
$x \mapsto |a(x)|\,|b(x)|$ is $\lambda_K$-integrable: it has the same absolute
value as $ab$, which is integrable by Step 1, so integrability transfers.

Step 3 (dominating the integral).  By hypothesis (2) the integrand $h$ is
dominated in absolute value by $C\,|a|\,|b|$, which is $\lambda_K$-integrable by
Step 2 and by $C \geq 0$.  The norm of an integral is at most the integral of
the norm of the integrand, and the integral is monotone on integrable
dominating pairs, so
\[
  \Bigl\| \int_K h \, d\lambda \Bigr\|
    \leq \int_K |h| \, d\lambda
    \leq \int_K C\,|a|\,|b| \, d\lambda .
\]
(If $h$ is not integrable the left side is $0$ and the chain holds trivially, as
noted in the statement.)

Step 4 (conclusion).  Pulling the nonnegative constant $C$ out of the last
integral by linearity and applying (CS),
\[
  \int_K C\,|a|\,|b| \, d\lambda
    = C \int_K |a|\,|b| \, d\lambda
    \leq C\,\|a\|_{L^2(K)}\,\|b\|_{L^2(K)} ,
\]
the final inequality being (CS) multiplied by $C \geq 0$, which preserves the
direction.  Chaining with Step 3 gives the assertion.
\end{proof}
